% GENERATED FILE. Edit canonical lessons, then run npm run generate:books.
% canonical-source-hash: fnv1a64:f68e20864c02b006
% canonical-lessons: JA-C446-kushami, JA-C446-hanamizu, JA-C446-samuke, JA-C446-ugai, JA-C446-taichou

\chapter{A Cold Coming On}
\label{ch:ja-a-cold-coming-on}
\hlchaptermodality{446}

\begin{chapteropening}
\textbf{By the end of this chapter:} \emph{I can say I am catching a cold.}

Complete the last lesson of chapter 446: I can say I am catching a cold.
\end{chapteropening}

\section[\texorpdfstring{\ja{くしゃみ} (kushami)}{くしゃみ (kushami)}]{\ja{くしゃみ} (kushami) — a sneeze}

\label{lesson:JA-C446-kushami}



\begin{quote}
Before the new one: say the Japanese for nausea, then the Japanese for a bruise.
\end{quote}

\subsection*{You'll want to know: \ja{くしゃみ}}
\textbf{\ja{くしゃみ}} — \emph{kushami} — \textquotedblleft{}a sneeze\textquotedblright{}.

\begin{quote}
\textbf{\ja{かぜで} \ja{くしゃみを} \ja{しました。}} — \emph{kaze de kushami o shimashita} — \textquotedblleft{}I sneezed because of my cold.\textquotedblright{}
\end{quote}

\begin{cousinweb}
\textbf{\ja{くしゃみ}} (\emph{kushami}) is a native Japanese word (wago), inherited rather than borrowed from Chinese; JMdict writes it in kana (JMdict).
\end{cousinweb}

\begin{grammarlens}[title={What you've built}]
One more word for describing what you lived and felt.
\end{grammarlens}

\subsection*{Your turn}
\begin{itemize}
\raggedright
  \item \emph{Say it:} \emph{kushami}
  \item \emph{Say it:} \emph{kushami}, once more
  \item \emph{Say it:} \emph{kaze de kushami o shimashita} — I sneezed because of my cold.
  \item \emph{Recall:} say the Japanese for nausea, then the Japanese for a bruise, then say \emph{kushami} again
\end{itemize}

\begin{tcolorbox}[breakable,colback=teal!4,colframe=teal!35!black,title={Before you move on}]
What does \ja{くしゃみ} mean? (\textquotedblleft{}A sneeze\textquotedblright{}.) Say it once more.
\end{tcolorbox}
\section[\texorpdfstring{\ja{はなみず} (hanamizu)}{はなみず (hanamizu)}]{\ja{はなみず} (hanamizu) — a runny nose}

\label{lesson:JA-C446-hanamizu}



\begin{quote}
Before the new one: say the Japanese for a bruise, then the Japanese for a sneeze.
\end{quote}

\subsection*{You'll want to know: \ja{はなみず}}
\textbf{\ja{はなみず}} — \emph{hanamizu} — \textquotedblleft{}a runny nose\textquotedblright{}.

\begin{quote}
\textbf{\ja{はなみずが} \ja{とまりません。}} — \emph{hanamizu ga tomarimasen} — \textquotedblleft{}My nose will not stop running.\textquotedblright{}
\end{quote}

\begin{cousinweb}
\textbf{\ja{はなみず}} (\emph{hanamizu}) is a native Japanese compound (wago), inherited rather than borrowed from Chinese: \emph{hana}, \textquotedblleft{}nose\textquotedblright{} and \emph{mizu}, \textquotedblleft{}water\textquotedblright{} (JMdict, KANJIDIC2).
\end{cousinweb}

\begin{grammarlens}[title={What you've built}]
One more word for describing what you lived and felt.
\end{grammarlens}

\subsection*{Your turn}
\begin{itemize}
\raggedright
  \item \emph{Say it:} \emph{hanamizu}
  \item \emph{Say it:} \emph{hanamizu}, once more
  \item \emph{Say it:} \emph{hanamizu ga tomarimasen} — My nose will not stop running.
  \item \emph{Recall:} say the Japanese for a bruise, then the Japanese for a sneeze, then say \emph{hanamizu} again
\end{itemize}

\begin{tcolorbox}[breakable,colback=teal!4,colframe=teal!35!black,title={Before you move on}]
What does \ja{はなみず} mean? (\textquotedblleft{}A runny nose\textquotedblright{}.) Say it once more.
\end{tcolorbox}
\section[\texorpdfstring{\ja{さむけ} (samuke)}{さむけ (samuke)}]{\ja{さむけ} (samuke) — a chill; the shivers}

\label{lesson:JA-C446-samuke}



\begin{quote}
Before the new one: say the Japanese for a sneeze, then the Japanese for a runny nose.
\end{quote}

\subsection*{You'll want to know: \ja{さむけ}}
\textbf{\ja{さむけ}} — \emph{samuke} — \textquotedblleft{}a chill; the shivers\textquotedblright{}.

\begin{quote}
\textbf{\ja{ねつが} \ja{あって、さむけが} \ja{します。}} — \emph{netsu ga atte, samuke ga shimasu} — \textquotedblleft{}I have a fever and the shivers.\textquotedblright{}
\end{quote}

\begin{cousinweb}
\textbf{\ja{さむけ}} (\emph{samuke}) is a mixed word: native \emph{samu}, \textquotedblleft{}cold\textquotedblright{}, the stem of \emph{samui}, joined to Sino-Japanese \emph{ke}, \textquotedblleft{}spirit, feeling\textquotedblright{} (Middle Chinese \emph{khj+jH}), a reading borrowed from Chinese (JMdict, KANJIDIC2; Middle Chinese in Baxter's transcription).
\end{cousinweb}

\begin{grammarlens}[title={What you've built}]
One more word for describing what you lived and felt.
\end{grammarlens}

\subsection*{Your turn}
\begin{itemize}
\raggedright
  \item \emph{Say it:} \emph{samuke}
  \item \emph{Say it:} \emph{samuke}, once more
  \item \emph{Say it:} \emph{netsu ga atte, samuke ga shimasu} — I have a fever and the shivers.
  \item \emph{Recall:} say the Japanese for a sneeze, then the Japanese for a runny nose, then say \emph{samuke} again
\end{itemize}

\begin{tcolorbox}[breakable,colback=teal!4,colframe=teal!35!black,title={Before you move on}]
What does \ja{さむけ} mean? (\textquotedblleft{}A chill; the shivers\textquotedblright{}.) Say it once more.
\end{tcolorbox}
\section[\texorpdfstring{\ja{うがい} (ugai)}{うがい (ugai)}]{\ja{うがい} (ugai) — gargling}

\label{lesson:JA-C446-ugai}



\begin{quote}
Before the new one: say the Japanese for a runny nose, then the Japanese for a chill.
\end{quote}

\subsection*{You'll want to know: \ja{うがい}}
\textbf{\ja{うがい}} — \emph{ugai} — \textquotedblleft{}gargling\textquotedblright{}.

\begin{quote}
\textbf{\ja{うちに} \ja{かえったら、うがいを} \ja{して} \ja{ください。}} — \emph{uchi ni kaettara, ugai o shite kudasai} — \textquotedblleft{}When you get home, please gargle.\textquotedblright{}
\end{quote}

\begin{cousinweb}
\textbf{\ja{うがい}} (\emph{ugai}) is a native Japanese word (wago), inherited rather than borrowed from Chinese; JMdict writes it in kana (JMdict).
\end{cousinweb}

\begin{grammarlens}[title={What you've built}]
One more word for describing what you lived and felt.
\end{grammarlens}

\subsection*{Your turn}
\begin{itemize}
\raggedright
  \item \emph{Say it:} \emph{ugai}
  \item \emph{Say it:} \emph{ugai}, once more
  \item \emph{Say it:} \emph{uchi ni kaettara, ugai o shite kudasai} — When you get home, please gargle.
  \item \emph{Recall:} say the Japanese for a runny nose, then the Japanese for a chill, then say \emph{ugai} again
\end{itemize}

\begin{tcolorbox}[breakable,colback=teal!4,colframe=teal!35!black,title={Before you move on}]
What does \ja{うがい} mean? (\textquotedblleft{}Gargling\textquotedblright{}.) Say it once more.
\end{tcolorbox}
\section[\texorpdfstring{\ja{たいちょう} (taichō)}{たいちょう (taichō)}]{\ja{たいちょう} (taichō) — physical condition; shape}

\label{lesson:JA-C446-taichou}



\begin{quote}
Before the new one: say the Japanese for a chill, then the Japanese for gargling.
\end{quote}

\subsection*{You'll want to know: \ja{たいちょう}}
\textbf{\ja{たいちょう}} — \emph{taichō} — \textquotedblleft{}physical condition; shape\textquotedblright{}.

\begin{quote}
\textbf{\ja{せんしゅうから} \ja{たいちょうが} \ja{よく} \ja{ありません。}} — \emph{senshū kara taichō ga yoku arimasen} — \textquotedblleft{}I have not been in good shape since last week.\textquotedblright{}
\end{quote}

\begin{cousinweb}
\textbf{\ja{たいちょう}} (\emph{taichō}) is Sino-Japanese (kango): two Chinese characters read with their on'yomi, readings borrowed from Chinese. They are \emph{tai}, \textquotedblleft{}body\textquotedblright{} (Middle Chinese \emph{thejX}) and \emph{chō}, \textquotedblleft{}harmonize\textquotedblright{} (\emph{dew}) (KANJIDIC2; Middle Chinese in Baxter's transcription).
\end{cousinweb}

\begin{grammarlens}[title={What you've built}]
One more word for describing what you lived and felt.
\end{grammarlens}

\subsection*{Your turn}
\begin{itemize}
\raggedright
  \item \emph{Say it:} \emph{taichō}
  \item \emph{Say it:} \emph{taichō}, once more
  \item \emph{Say it:} \emph{senshū kara taichō ga yoku arimasen} — I have not been in good shape since last week.
  \item \emph{Recall:} say the Japanese for a chill, then the Japanese for gargling, then say \emph{taichō} again
\end{itemize}

\begin{tcolorbox}[breakable,colback=teal!4,colframe=teal!35!black,title={Before you move on}]
What does \ja{たいちょう} mean? (\textquotedblleft{}Physical condition; shape\textquotedblright{}.) Say it once more.
\end{tcolorbox}
